% The output of the desktop editor in a document that has its own \pkey and \set.
\documentclass{article}
\usepackage{erdiagram}
\newcommand\set[1]{\left\{#1\right\}}
\newcommand\pkey[1]{MINE(#1)}
\newcommand\relation{MINE}
\begin{document}
Before: \pkey{a}, $\set{x}$, \relation, \emph{emphasis stays italic}.
\begin{relationalschema}
  \relation{Person}{\pkey{PID}{integer}, \col{Name}{string}, \fkey{Raum-Id}{}}
  \relation{Hund}{\pfkey{PID}{integer}, \pkey{Nr}{}, \col{Vorgänger}{}}
  \relation{Very\_long}{\pkey{aaaaaaaaaaaaaaaaaaaa}{integer}, \col{bbbbbbbbbbbbbbbbbbbbbbbb}{string},
    \col{cccccccccccccccccccc}{string}, \col{dddddddddddddddddddddd}{string}, \col{eeeeeeeeeee}{string}}
\end{relationalschema}
After: \pkey{a}, $\set{x}$, \relation.
\end{document}
